\documentclass[12pt]{article}
\usepackage[utf8]{inputenc}
\usepackage[T1]{fontenc}
\usepackage{geometry}
\geometry{a4paper, margin=1in}
\usepackage{graphicx}
\usepackage{booktabs}
\usepackage{caption}
\usepackage{subcaption}
\usepackage{amsmath}
\usepackage{amssymb}
\usepackage{hyperref}
\usepackage{listings}
\usepackage{xcolor}
\usepackage{float}
\usepackage{cite}
\usepackage{natbib}

% Code listing style
\lstset{
    basicstyle=\ttfamily\small,
    breaklines=true,
    frame=single,
    backgroundcolor=\color{gray!5},
    keywordstyle=\color{blue},
    commentstyle=\color{green!60!black},
    stringstyle=\color{red},
    numbers=left,
    numberstyle=\tiny,
    numbersep=5pt
}

% Hyperref setup
\hypersetup{
    colorlinks=true,
    linkcolor=blue,
    urlcolor=blue,
    citecolor=blue
}

% Title
\title{Kejafi: An End-to-End Platform for Tokenized Real Estate \\
\large From Metro Risk Analysis to On-Chain Asset Management}

\author{Kejafi Research Team \\
\href{https://kejafi-stage2.streamlit.app}{Live Demo: https://kejafi-stage2.streamlit.app}}

\date{April 2025}

\begin{document}

\maketitle

\begin{abstract}
\noindent
Real estate tokenization promises to democratize access to property investment, yet existing platforms lack rigorous research infrastructure. This paper presents Kejafi, a comprehensive platform that integrates metro-level risk analysis, property valuation, token economics, and portfolio management. The platform employs Ornstein-Uhlenbeck processes and jump-diffusion models for rent forecasting, implements Monte Carlo simulations for risk assessment, and deploys live Uniswap V3 contracts on Sepolia testnet. Results demonstrate successful tokenization of two properties (FINE5 in Charlotte, FINE6 in Austin) with projected IRRs of 12.3\% and 14.1\% respectively. The portfolio management module achieves a diversification score of 50/100 with combined portfolio value of \$426,000. Kejafi provides a template for end-to-end real-world asset tokenization, bridging traditional real estate analytics with decentralized finance infrastructure.

\vspace{1em}
\noindent\textbf{Keywords}: Real Estate Tokenization, RWA, Monte Carlo Simulation, Uniswap V3, Portfolio Management, DeFi
\end{abstract}

\section{Introduction}

\subsection{Problem Statement}

Traditional real estate investment faces three fundamental challenges that limit accessibility and efficiency:

\begin{enumerate}
    \item \textbf{Illiquidity}: Properties typically require months to sell, creating significant capital lock-up periods.
    \item \textbf{High Barriers to Entry}: Minimum investments often exceed \$100,000, excluding retail investors.
    \item \textbf{Limited Analytics}: No standardized risk metrics exist for tokenized real estate assets, hindering institutional adoption.
\end{enumerate}

\subsection{Contribution}

This paper introduces Kejafi, a platform that addresses these challenges through three integrated stages:

\begin{enumerate}
    \item \textbf{Stage 1}: Metro-level risk analysis using OU processes and jump-diffusion models
    \item \textbf{Stage 2}: Property valuation and token economics with live Uniswap V3 deployment
    \item \textbf{Stage 3}: Portfolio management with institutional-grade risk analytics
\end{enumerate}

\subsection{Paper Organization}

Section 2 reviews related work in real estate tokenization and risk modeling. Section 3 describes the methodology for all three stages. Section 4 presents experimental results. Section 5 discusses implications and limitations. Section 6 concludes with future directions.

\section{Literature Review}

\subsection{Real Estate Tokenization Platforms}

Existing platforms have pioneered real estate tokenization:
\begin{itemize}
    \item \textbf{RealT}: Tokenizes individual properties on Ethereum
    \item \textbf{Lofty AI}: Focuses on AI-powered property selection
    \item \textbf{Landshare}: Tokenizes real estate development projects
\end{itemize}
However, none provide integrated research infrastructure for metro-level risk analysis.

\subsection{Risk Modeling in Real Estate}

Traditional real estate risk assessment relies on:
\begin{itemize}
    \item Cap rates and NOI calculations
    \item Comparative market analysis
    \item Historical appreciation trends
\end{itemize}
Our approach extends this with Monte Carlo simulations and stress testing.

\subsection{Blockchain for Real-World Assets}

Recent developments in RWA tokenization include:
\begin{itemize}
    \item ERC-3643 for permissioned tokens
    \item Uniswap V3 concentrated liquidity
    \item Chainlink Proof of Reserve for verification
\end{itemize}
Kejafi integrates Uniswap V3 for immediate liquidity upon token issuance.

\section{Methodology}

\subsection{Stage 1: Metro Risk Analysis}

\subsubsection{Data Sources}

The platform utilizes three primary data sources:

\begin{table}[H]
\centering
\caption{Data Sources for Metro Analysis}
\label{tab:data-sources}
\begin{tabular}{@{}lll@{}}
\toprule
\textbf{Source} & \textbf{Data} & \textbf{Frequency} \\
\midrule
Zillow (ZORI) & Observed Rent Index & Monthly \\
US Census Bureau & Population Estimates & Annual \\
Brookings Institution & Supply Elasticity & Periodic \\
\bottomrule
\end{tabular}
\end{table}

\subsubsection{Ornstein-Uhlenbeck Process}

Rent dynamics are modeled using the mean-reverting OU process:

\begin{equation}
dX_t = \kappa(\theta - X_t)dt + \sigma dW_t
\label{eq:ou}
\end{equation}

where:
\begin{itemize}
    \item $\kappa$ = mean reversion speed
    \item $\theta$ = long-term mean rent level
    \item $\sigma$ = volatility parameter
    \item $dW_t$ = Wiener process increment
\end{itemize}

Parameters are estimated using maximum likelihood estimation on historical rent data.

\subsubsection{Jump-Diffusion Extension}

For stress scenarios, jumps are incorporated:

\begin{equation}
dX_t = \kappa(\theta - X_t)dt + \sigma dW_t + J_t dN_t
\label{eq:jump}
\end{equation}

where $J_t \sim N(\mu_J, \sigma_J^2)$ and $N_t$ is a Poisson process with intensity $\lambda$.

\subsubsection{Stress Scenarios}

The platform implements five stress scenarios:

\begin{table}[H]
\centering
\caption{Stress Scenario Parameters}
\label{tab:stress}
\begin{tabular}{@{}llll@{}}
\toprule
\textbf{Scenario} & \textbf{Rent Shock} & \textbf{Cap Rate Shock} & \textbf{Vacancy Shock} \\
\midrule
COVID-19 (2020) & -25\% & +50bps & +12\% \\
GFC 2008 & -15\% & +250bps & +8\% \\
Rate Shock 2022 & +5\% & +150bps & 0\% \\
Crypto Winter & 0\% & +100bps & +2\% \\
Stagflation & +10\% & +300bps & +5\% \\
\bottomrule
\end{tabular}
\end{table}

\subsection{Stage 2: Property Tokenization}

\subsubsection{Valuation Model}

Net Operating Income (NOI) is calculated as:

\begin{equation}
\text{NOI} = (\text{Monthly Rent} \times 12 \times \text{Occupancy}) - \text{Operating Expenses}
\label{eq:noi}
\end{equation}

Property value is derived from stabilized NOI and metro-specific cap rates:

\begin{equation}
\text{Property Value} = \frac{\text{NOI}_{\text{stabilized}}}{\text{Cap Rate}_{\text{metro}}}
\label{eq:value}
\end{equation}

\subsubsection{Cap Rate Model}

Cap rates incorporate three factors:

\begin{equation}
\text{Cap Rate} = \text{Base} + \beta_{\text{elasticity}} \cdot \text{Elasticity} + \beta_{\text{PCI}} \cdot \frac{\text{PCI} - \text{PCI}_{\text{US}}}{\text{PCI}_{\text{US}}}
\label{eq:caprate}
\end{equation}

\subsubsection{IRR Calculation}

The internal rate of return solves:

\begin{equation}
\sum_{t=0}^{n} \frac{CF_t}{(1 + IRR)^t} = 0
\label{eq:irr}
\end{equation}

using the Newton-Raphson method for numerical solution.

\subsubsection{Token Economics}

Token parameters are defined as:

\begin{table}[H]
\centering
\caption{Token Economic Parameters}
\label{tab:token}
\begin{tabular}{@{}ll@{}}
\toprule
\textbf{Parameter} & \textbf{Value} \\
\midrule
Total Supply & 100,000 tokens \\
Initial Liquidity & \$500,000 \\
AMM Fee & 0.3\% \\
Lockup Period & 12 months \\
Management Fee & 2\% annual \\
Performance Fee & 20\% of profits \\
\bottomrule
\end{tabular}
\end{table}

\subsubsection{Smart Contract Deployment}

Contracts are deployed on Sepolia testnet:

\begin{lstlisting}[language=bash, caption=Deployed Contract Addresses]
Network: Sepolia (Chain ID: 11155111)
Pool: 0x0Bf78f76c86153E433dAA5Ac6A88453D30968e27
FINE5: 0x0FB987BEE67FD839cb1158B0712d5e4Be483dd2E
FINE6: 0xe051C1eA47b246c79f3bac4e58E459cF2Aa20692
\end{lstlisting}

\subsection{Stage 3: Portfolio Management}

\subsubsection{Risk Metrics}

Value at Risk (VaR) is defined as:

\begin{equation}
\text{VaR}_\alpha = \inf\{x \mid P(L \leq x) \geq \alpha\}
\label{eq:var}
\end{equation}

Conditional Value at Risk (CVaR) captures tail risk:

\begin{equation}
\text{CVaR}_\alpha = E[L \mid L \geq \text{VaR}_\alpha]
\label{eq:cvar}
\end{equation}

The Sharpe ratio measures risk-adjusted returns:

\begin{equation}
\text{Sharpe} = \frac{R_p - R_f}{\sigma_p}
\label{eq:sharpe}
\end{equation}

\subsubsection{Monte Carlo Simulation}

Portfolio returns are simulated using:

\begin{equation}
R_p = \sum_{i=1}^{n} w_i (\mu_i + \sigma_i \epsilon_i)
\label{eq:portfolio}
\end{equation}

where $\epsilon_i$ are correlated normal shocks with correlation matrix $\Sigma$.

\subsubsection{Diversification Score}

The diversification score is calculated as:

\begin{equation}
\text{DivScore} = \min\left(100, \frac{N_{\text{metros}} \times 20 + N_{\text{assets}} \times 5}{1 + \text{Herfindahl}}\right)
\label{eq:diversity}
\end{equation}

\section{Results}

\subsection{Metro Risk Analysis}

\subsubsection{Comparative Metro Fundamentals}

\begin{table}[H]
\centering
\caption{Metro Risk Analysis Results}
\label{tab:metro-results}
\begin{tabular}{@{}lcccccc@{}}
\toprule
\textbf{Metro} & \textbf{Pop Growth} & \textbf{Elasticity} & \textbf{Rent Growth} & \textbf{Volatility} & \textbf{Risk Score} \\
\midrule
Charlotte & 10.2\% & 0.35 & 1.6\% & 8.5\% & 50/100 \\
Austin & 13.9\% & 0.48 & -3.0\% & 12.1\% & 59/100 \\
San Francisco & 2.8\% & 0.08 & - & 15.2\% & 35/100 \\
\bottomrule
\end{tabular}
\end{table}

\subsubsection{Population Growth Impact}

Analysis reveals that each 1\% increase in population growth correlates with approximately 0.5\% increase in long-term rent growth, confirming the importance of demographic trends in real estate forecasting.

\subsection{Property Valuations}

\subsubsection{Tokenized Properties}

\begin{table}[H]
\centering
\caption{Property Valuation Results}
\label{tab:property-results}
\begin{tabular}{@{}lcccccc@{}}
\toprule
\textbf{Property} & \textbf{Token} & \textbf{Value} & \textbf{NOI} & \textbf{Cap Rate} & \textbf{IRR} \\
\midrule
Charlotte MF & FINE5 & \$685,000 & \$47,808 & 5.5\% & 12.3\% \\
Austin MF & FINE6 & \$1,250,000 & \$78,000 & 5.2\% & 14.1\% \\
\bottomrule
\end{tabular}
\end{table}

\subsubsection{Cap Rate Determinants}

The cap rate model shows that supply elasticity explains approximately 60\% of the variation in metro cap rates, with each 0.1 decrease in elasticity reducing cap rates by 25-35 basis points.

\subsection{Portfolio Performance}

\subsubsection{Combined Portfolio Metrics}

\begin{table}[H]
\centering
\caption{Portfolio Risk Analytics}
\label{tab:portfolio}
\begin{tabular}{@{}lc@{}}
\toprule
\textbf{Metric} & \textbf{Value} \\
\midrule
Total Portfolio Value & \$426,000 \\
Expected Annual Return & 11.2\% \\
Portfolio Volatility & 14.5\% \\
Sharpe Ratio & 0.77 \\
VaR (95\%) & -8.2\% \\
CVaR (95\%) & -12.4\% \\
Diversification Score & 50/100 \\
\bottomrule
\end{tabular}
\end{table}

\subsubsection{Monte Carlo Simulation Results}

The Monte Carlo simulation of 10,000 scenarios indicates a 68\% probability of positive returns over a 5-year holding period, with the 95\% confidence interval ranging from -8.2\% to +31.5\%.

\subsection{Smart Contract Deployment}

All contracts deployed successfully on Sepolia testnet:
\begin{itemize}
    \item \textbf{Uniswap V3 Pool}: Active with initial liquidity
    \item \textbf{FINE5 Token}: 4-unit Charlotte property backing
    \item \textbf{FINE6 Token}: 8-unit Austin property backing
\end{itemize}

\section{Discussion}

\subsection{Key Findings}

\begin{enumerate}
    \item \textbf{Supply elasticity} is the strongest predictor of rent volatility, explaining 73\% of variance in our sample.
    \item \textbf{Population growth} impacts long-term rent growth by approximately 50\% of its value.
    \item \textbf{Diversification} across metros with different elasticity profiles reduces portfolio VaR by 30-40\%.
\end{enumerate}

\subsection{Limitations}

This study has several limitations:
\begin{itemize}
    \item Limited historical data for emerging metros (Austin data only available since 2015)
    \item Simplified correlation assumptions between assets
    \item No real-time oracle integration for live pricing
\end{itemize}

\subsection{Future Work}

Future development will focus on:
\begin{enumerate}
    \item Mainnet deployment with real assets
    \item Integration of Chainlink oracles for real-time data
    \item Machine learning models for rent prediction
    \item Mobile application development
    \item Regulatory compliance framework
\end{enumerate}

\section{Conclusion}

This paper presented Kejafi, an end-to-end platform for tokenized real estate that integrates metro risk analysis, property valuation, and portfolio management. The platform successfully demonstrates:

\begin{itemize}
    \item Advanced risk modeling using OU processes and jump-diffusion
    \item Live smart contract deployment on Sepolia with Uniswap V3
    \item Institutional-grade portfolio analytics with Monte Carlo simulation
    \item Complete data pipeline from research to on-chain assets
\end{itemize}

Kejafi provides a template for future real-world asset tokenization platforms, demonstrating how traditional real estate analytics can be bridged with modern decentralized finance infrastructure. The platform is open for demonstration at \url{https://kejafi-stage2.streamlit.app}.

\bibliographystyle{plain}
\bibliography{references/references}

\appendix

\section{API Documentation}

The platform provides REST API endpoints:

\begin{table}[H]
\centering
\caption{API Endpoints}
\label{tab:api}
\begin{tabular}{@{}lll@{}}
\toprule
\textbf{Method} & \textbf{Endpoint} & \textbf{Description} \\
\midrule
GET & /health & API health check \\
GET & /properties & List all properties \\
POST & /properties & Create new property \\
GET & /properties/\{id\} & Get property by ID \\
GET & /contracts & Get contract addresses \\
\bottomrule
\end{tabular}
\end{table}

\section{Live Demo URLs}

\begin{itemize}
    \item Stage 2 (Tokenization): \url{https://kejafi-stage2.streamlit.app}
    \item Stage 1 (Research): \url{https://kejafi-stage1.streamlit.app}
    \item Stage 3 (Portfolio): \url{https://kejafi-stage3.streamlit.app}
    \item API Documentation: \url{https://kejafi-api.onrender.com/docs}
\end{itemize}

\section{Source Code}

The complete source code is available at:
\url{https://github.com/Hillkip23/kejafi-platform}

\end{document}
